\section{集合族的并、交、积}

\begin{enumerate}
\def\labelenumi{\arabic{enumi}.}
\tightlist
\item
  在本节中，我们考虑集合E的子集的一个族 \((\mathbf{X}_{i})_{i\in\mathbf{I}}\) ，其中指标集I是任意的；属于该族的E的子集所构成的集合将记为 \(\mathfrak{F}\) （它因此是 \(\mathfrak{P}(\mathbf{E})\) 的一个子集）.
\end{enumerate}

如果 I 是有限的，那么考虑族 \((\mathbf{X}_{t})\) 归结为E的若干子集的问题，这些子集可能相同也可能不同，其数目等于I的元素个数。例如，E 的任意三个子集 \(X_{1}, X_{2}, X_{3}\) 构成 E 的一个子集族，这里的指标集由数字 1、2、3 组成。

\begin{enumerate}
\def\labelenumi{\arabic{enumi}.}
\setcounter{enumi}{1}
\tightlist
\item
  让 \(J\) 是 \(I\) 的任意子集，并考虑所有元素 \(x\) 组成的集合，这些元素具有性质“存在 \(\iota \in J\) 使得 \(x \in X_\iota\)''。这个集合被称为集合族 \((X_\iota)_{\iota \in J}\) 的并集，并记作 \(\bigcup_{\iota\in J} X_\iota\) .
\end{enumerate}

我们也可以将这一表述定义如下：对于映射 \(\iota \to \mathbf{X}_{\iota}\) （从 \(\mathbf{I}\) 到 \(\mathfrak{P}(\mathbf{E})\) ）对应着一个良定义的子集 \(\mathbf{C}\) （ \(\mathbf{I} \times \mathbf{E}\) 的）使得 \(\mathbf{X}_{\iota} = \mathbf{C}(\iota)\) ( \(\S 3\) ，第7号），并且我们有 \(\bigcup_{\iota \in J} \mathbf{X}_{\iota} = \mathbf{C}(\mathbf{J})\) .

特别是，

\[
\bigcup_ {t \in \varnothing} \mathrm{X} _ {t} = \mathrm{C} (\varnothing) = \varnothing .\tag{33}
\]

当 \(J = I\) ，我们常写作 \(\bigcup_{i} X_i\) ，或简称为 \(\bigcup X_i\) ，代替 \(\bigcup_{i \in I} X_i\) .

并集 \(\bigcup_{i} X_{i}\) 仅依赖于集合 \(\mathfrak{F}\) ；换言之，对于对应于同一子集 \(\mathfrak{F}\) （ \(\mathfrak{P}(E)\) 的）的两个族，它是相同的。特别地，它等于由 \(\mathfrak{F}\) 到 \(\mathfrak{P}(E)\) 的典范映射所定义的族之并，因此我们可以写成 \(\bigcup_{X \in \mathfrak{F}} X\) ，它也称为属于 \(\mathfrak{F}\) 的集合之并.

当 \(\mathbf{I}\) 是一个其元素被明确指定的集合，例如数1、2、3，我们有

\[
\bigcup_ {i \in I} X _ {i} = X _ {1} \cup X _ {2} \cup X _ {3},
\]

这证明了在一般情况下赋予该集合的名称“并”是合理的。 \(\bigcup_{i}X_{i}\) .

\begin{enumerate}
\def\labelenumi{\arabic{enumi}.}
\setcounter{enumi}{2}
\tightlist
\item
  对于任意 \(J \subset I\) 我们拥有
\end{enumerate}

\[
\bigcup_ {i \in J} X _ {i} \subset \bigcup_ {i \in I} X _ {i}.
\]

特别地，对于所有 \(x \in I\) ，我们有

\[
\mathrm{X} _ {x} \subset \bigcup_ {i \in I} \mathrm{X} _ {i}.
\]

反之，如果 Y 是 E 的子集且满足 \(X_{i} \subset Y\) 对于所有 \(i \in I\) ，那么

\[
\bigcup_ {i} \mathrm{X} _ {i} \subset \mathrm{Y}.
\]

更一般地，如果 \((\mathbf{Y}_{i})\) 是E的另一族子集，由同一指标集I索引，并且如果 \(X_{i} \subset Y_{i}\) 对于所有 \(i \in I\) ，那么

\[
\bigcup_ {i} X _ {i} \subset \bigcup_ {i} Y _ {i}.
\]

设 F 为另一个集合，并设 \(\mathbf{X} \to \mathbf{K}(\mathbf{X})\) 为从 \(\mathfrak{P}(\mathbf{E})\) 到 \(\mathfrak{P}(\mathbf{F})\) 的映射，它由 \(E \times F\) 的子集 K 定义。于是我们有

\[
\mathrm{K} \left(\bigcup_ {i \in J} \mathrm{X} _ {i}\right) = \bigcup_ {i \in J} \mathrm{K} (\mathrm{X} _ {i}).\tag{34}
\]

现在令 \(\mathbf{L}\) 为另一个指标集，且 \((J_{\lambda})_{\lambda \in \mathbf{L}}\) 为 \(\mathbf{I}\) 。然后

\[
\bigcup_ {\iota \in \bigcup_ {\lambda \in L} J _ {\lambda}} X _ {\iota} = \bigcup_ {\lambda \in L} \left(\bigcup_ {\iota \in J _ {\lambda}} X _ {\iota}\right).\tag{35}
\]

的子集族。这是并集的一般结合律公式。当 I 和 L 是元素被明确指定的集合时，我们得到的这些关系已经给出（见第 1 节，第 14 号）。如果仅 L 满足此条件，且例如由数 1、2 组成，则我们有

\[
\bigcup_ {i \in J _ {1} \cup J _ {2}} X _ {i} = \left(\bigcup_ {i \in J _ {1}} X _ {i}\right) \cup \left(\bigcup_ {i \in J _ {2}} X _ {i}\right).\tag{36}
\]

令 \((\mathbf{X}_i)_{i\in I}\) 与 \((\mathbf{Y}_x)_{x\in K}\) 为E的任意两个子集族。则我们有

\[
\left(\bigcup_ {i \in I} X _ {i}\right) \cap \left(\bigcup_ {x \in K} Y _ {x}\right) = \bigcup_ {(i, x) \in I \times K} (X _ {i} \cap Y _ {x}),\tag{37}
\]

分配律公式，其中包含(10)中第二个公式作为特例。

如果 \((\mathbf{X}_t)_{i\in \mathbf{I}}\) 是以下集合的子集族 \(\mathbf{E}\) ，并且如果 \((\mathbf{Y}_{\mathbf{x}})_{\mathbf{x}\in \mathbf{K}}\) 是以下集合的子集族 \(\mathbf{F}\) ，那么

\[
\left(\bigcup_ {i \in I} X _ {i}\right) \times \left(\bigcup_ {x \in K} Y _ {x}\right) = \bigcup_ {(i, x) \in I \times K} (X _ {i} \times Y _ {x}).\tag{38}
\]

\begin{enumerate}
\def\labelenumi{\arabic{enumi}.}
\setcounter{enumi}{3}
\tightlist
\item
  一个族 \((\mathbf{X}_i)_{i\in I}\) 子集的 \(E\) 是子集的一个覆盖 \(A\) 的 \(E\) ，或覆盖 \(A\) ，如果
\end{enumerate}

\[
\mathrm{A} \subset \bigcup_ {i \in \mathrm{I}} \mathrm{X} _ {i}.
\]

特别地，如果 \((\mathbf{X}_{t})\) 是E的一个覆盖，我们有

\[
\bigcup_ {i \in I} X _ {i} = E.
\]

一个划分， \(\mathbf{E}\) 是一个覆盖 \((\mathbf{X}_t)\) 的 \(\mathbf{E}\) 使得

\begin{enumerate}
\def\labelenumi{(\alph{enumi})}
\item
  \(\mathbf{X}_{\iota} \neq \varnothing\) 对于所有 \(\iota \in \mathbf{I}\) ;
\item
  \(X_{t} \cap X_{x} = \varnothing\) 对于I中每一对不同指标t、x。（这第二个条件可以表述为： \(X_{t}\) 是两两不相交的。
\end{enumerate}

这些条件蕴含 \(\iota \to X_{\iota}\) 是 I 到该集合的双射 \(\mathfrak{F}\) 划分的子集。因此，若 \(\mathfrak{F}\) 给定，则该族在索引集的一一对应意义下被确定。特别地，一个划分既可被视为子集的集合，也可被视为子集的族。

\begin{enumerate}
\def\labelenumi{\arabic{enumi}.}
\setcounter{enumi}{4}
\tightlist
\item
  令 \((\mathbf{X}_t)_{i \in I}\) 为集合 \(E\) 的任意非空子集族。在乘积 \(I \times E\) 中，令 \(\mathbf{X}_t'\) 表示子集 \(\{\iota\} \times \mathbf{X}_t\) （对于每一个 \(\iota \in I\) ）。集合
\end{enumerate}

\[
S = \bigcup_ {i \in I} X _ {i} ^ {\prime}
\]

被称为族 \((\mathbf{X}_i)_{i\in I}\) 的和。显然，族 \((\mathbf{X}_i')_{i\in I}\) 构成 \(S\) 的一个划分，且对每个 \(i\in I\) ，映射 \(x_i\to (i,x_i)\) 是 \(\mathbf{X}_i\) 到 \(\mathbf{X}_i'\) 上的一个双射。滥用语言，任何与 \(S\) 一一对应的集合通常被称为族 \((\mathbf{X}_i)_{i\in I}\) 的和，并且 \(\mathbf{X}_i\) 通常被等同于这个集合中与之对应的子集。

两个非空集合E和F的和通常被称为通过将集合F附加到E上而得到的。

\begin{enumerate}
\def\labelenumi{\arabic{enumi}.}
\setcounter{enumi}{5}
\tightlist
\item
  利用第2节的记号，满足性质“对所有 \(\iota\in J\) , \(x\in X_{\iota}\) ”的E的元素x的集合称为集合族 \((\mathbf{X}_i)_{i \in J}\) 的交集，并记作 \(\bigcap_{i \in J} \mathbf{X}_i\) ；当 \(J = I\) ，我们常写作 \(\bigcap_{i} \mathbf{X}_i\) ，或简称为 \(\bigcap \mathbf{X}_i\) ，而不是 \(\bigcap_{i \in I} \mathbf{X}_i\) .
\end{enumerate}

我们有

\[
\mathbb {C} \left(\bigcup_ {i \in J} X _ {i}\right) = \bigcap_ {i \in J} \left(\mathbb {C} X _ {i}\right).\tag{39}
\]

特别地，如果 \(\mathrm{J} = \emptyset\) ,

\[
\bigcap_ {i \in \emptyset} X _ {i} = E.\tag{40}
\]

交集 \(\bigcap_{i} X_i\) 仅依赖于集合 \(\mathfrak{F}\) ，并且可以写成 \(\bigcap_{X \in \mathfrak{F}} X\) 。例如，如果 \(I\) 由数字1、2、3组成，我们有

\[
\bigcap_ {i} \mathrm{X} _ {i} = \mathrm{X} _ {1} \cap \mathrm{X} _ {2} \cap \mathrm{X} _ {3}.
\]

\begin{enumerate}
\def\labelenumi{\arabic{enumi}.}
\setcounter{enumi}{6}
\tightlist
\item
  公式(39)使我们能够推广对偶规则。如果E的子集A是由其他子集X、Y、Z及E的子集族 \((\mathbf{X}_{t}), (\mathbf{Y}_{x}), (\mathbf{Z}_{\lambda})\) 仅通过（以任意顺序）应用运算 \(\mathbb C, \cup, \cap, \bigcup, \bigcap\) 而得到，那么我们可以如下得到补集 \(\mathbb C A\) ：将子集 X、Y、Z， \(X_{t}, Y_{x}, Z_{\lambda}\) 替换为它们的补集，并将运算 \(\cup, \cap, \bigcup, \bigcap\) 分别替换为 \(\cap, \cup, \bigcap, \bigcup\) ，运算顺序保持不变；当然，当交与并的运算作用于写在符号 \(\bigcup\) 与 \(\bigcap\) 下方的指标集时，这些运算不得改变。
\end{enumerate}

如第1节第15号所述，我们定义关系 \(\mathbf{A} = \mathbf{B}\) 或 \(\mathbf{A} \subset \mathbf{B}\) （其中 \(\mathbf{A}\) 与 \(\mathbf{B}\) 是 \(\mathbf{E}\) 的上述形式的子集）的对偶。

\begin{enumerate}
\def\labelenumi{\arabic{enumi}.}
\setcounter{enumi}{7}
\tightlist
\item
  对于所有 \(J \subset I\) 我们拥有
\end{enumerate}

\[
\bigcap_ {i \in I} X _ {i} \subset \bigcap_ {i \in J} X _ {i}.
\]

特别地，对于所有 \(x \in I\) 我们拥有

\[
\bigcap_ {i} \mathrm{X} _ {i} \subset \mathrm{X} _ {x}.
\]

反之，如果 \(\mathbf{Y} \subset \mathbf{X}_i\) 对于所有 \(\iota \in I\) ，那么

\[
\mathrm{Y} \subset \bigcap_ {i} \mathrm{X} _ {i}.
\]

更一般地，如果 \((\mathbf{Y}_{t})\) 是E的另一族子集，由同一指标集I索引，并且如果 \(X_{t} \subset Y_{t}\) 对于所有 \(t \in I\) ，那么

\[
\bigcap_ {i} X _ {i} \subset \bigcap_ {i} Y _ {i}.
\]

这些集合的并集 \(X_{t}\) 是所有满足以下条件的集合Y的交集 \(X_{t} \subset Y\) 对于所有 \(t \in I\) 。这些集合的交集 \(X_{t}\) 是所有满足以下条件的Z的并集 \(Z \subset X_{t}\) 对于所有 \(t \in I\) .

以下公式分别是（35）和（37）的对偶式：

\begin{enumerate}
\def\labelenumi{(\arabic{enumi})}
\setcounter{enumi}{40}
\tightlist
\item
\end{enumerate}

\[
\bigcap_ {i \in \bigcup_ {\lambda \in L} J _ {\lambda}} X _ {i} = \bigcap_ {\lambda \in L} \left(\bigcap_ {i \in J _ {\lambda}} X _ {i}\right) \quad (\text { 结合律 });\tag{42}
\]

\[
\left(\bigcap_ {t \in I} X _ {t}\right) \cup \left(\bigcap_ {x \in K} Y _ {x}\right) = \bigcap_ {(t, x) \in I \times K} (X _ {t} \cup Y _ {x}) \quad (\text { 分配律 }).
\]

如果 \((\mathbf{X}_i)_{i\in I}\) 是以下集合的子集族 \(\mathbf{E}\) ，以及 \((\mathbf{Y}_x)_{x\in \mathbb{K}}\) 一族子集 \(\mathbf{F}\) ，那么

\[
\left(\bigcap_ {i \in I} X _ {i}\right) \times \left(\bigcap_ {x \in K} Y _ {x}\right) = \bigcap_ {(i, x) \in I \times K} (X _ {i} \times Y _ {x}).\tag{43}
\]

此外，如果 \((\mathbf{X}_{t})\) 并且 \((\mathbf{Y}_{t})\) 分别是E和F的子集族，由同一指标集I索引，则

\[
\left(\bigcap_ {i \in I} X _ {i}\right) \times \left(\bigcap_ {i \in I} Y _ {i}\right) = \bigcap_ {i \in I} (X _ {i} \times Y _ {i}).\tag{44}
\]

公式(34)没有对偶；一般而言，我们只能说

\[
\mathrm{K} \left(\bigcap_ {i \in J} X _ {i}\right) \subset \bigcap_ {i \in J} \mathrm{K} (X _ {i}).\tag{45}
\]

对于所有族，我们在(45)中都有等式成立 \((\mathbf{X}_{t})\) 仅当 \(\mathbf{X} \to \mathbf{K}(\mathbf{X})\) 是从F的子集到E的映射的逆扩张。因此，如果f是从F到E的映射，我们有（推广(14)）

\[
\bar {f} ^ {1} \left(\bigcap_ {i \in J} X _ {i}\right) = \bigcap_ {i \in J} \bar {f} ^ {1} (X _ {i}).\tag{46}
\]

\begin{enumerate}
\def\labelenumi{\arabic{enumi}.}
\setcounter{enumi}{8}
\tightlist
\item
  设E为任意集合，I为任意指标集。以 I 为指标集的、E 的元素的所有族 \((x_{i})_{i\in I}\) 所构成的集合记作 \(E^{I}\) ，而从 E 到 \(E^{I}\) 的过渡运算称为幂运算。 \(E^{I}\) 因此与从I到E的所有映射的集合一一对应（该集合因此通常也记作 \(E^{I}\) ，这是滥用语言的说法），同时也与 \(\mathfrak{P}(\mathbf{I} \times \mathbf{E})\) 的某个子集一一对应，只需考虑这些映射的图像。对应于集合 I 的子集 J 的那些集合 \(E^{J}\) 因此都可以被视为同一个集合的子集，而该集合与 \(\mathfrak{P}(\mathbf{I} \times \mathbf{E})\) 的某个子集一一对应.
\end{enumerate}

现在令 \((\mathbf{X}_i)_{i\in I}\) 是 E 的子集族，由同一个集合 I 索引，并设 J 是 I 的任意子集。性质“对所有 \(i\in J\) , \(x_i\in X_i\)''（族 \((x_i)_{i\in J}\) 的）定义了 \(\mathbf{E}^{\mathbf{J}}\) 的一个子集，称为集合族 \((\mathbf{X}_i)_{i\in J}\) 的乘积，并记作 \(\prod_{i\in J}\mathbf{X}_i\) （或简记为 \(\prod_{i}\mathbf{X}_i\) 当 \(J = I\) ）。 \(\mathbf{X}_i\) 被称为乘积的因子。注意， \(\prod_{i\in \emptyset}\mathbf{X}_i\) 是一个由单个元素组成的集合（对应于 \(\mathbf{I}\times \mathbf{E}\) 的空子集）。如果 \(\mathbf{X}_i = \mathbf{E}\) （对所有 \(i\in J\) ），那么我们有

\[
\prod_ {i \in J} X _ {i} = E ^ {J}.
\]

如果，例如，I 由三个数 1、2、3 组成，那么 \(\prod_{i\in J}X_i\) 与集合 \(X_1\times X_2\times X_3\) 一一对应。10. 如果 \(R\{x,y\}\) 是集合 E 的一个一般元素 \(x\) 与集合 F 的一个一般元素 \(y\) 之间的关系，则以下命题等价：

“对每个 x，存在 y 使得 R\{x,y\}''

并且

“存在一个从 E 到 F 的映射 f，使得 R\{x, f(x)\} 对于所有 x''。

这一等价性的断言被称为选择公理（或策梅洛公理）。我们有时会指出某个定理的证明是否依赖于它。

选择公理等价于以下命题：“如果对每个 \(\iota\in I\) ，我们有 \(X_{\iota}\neq\phi\) ，那么 \(\prod_{\iota\in I}X_{\iota}\neq\phi\) ''.

\begin{enumerate}
\def\labelenumi{\arabic{enumi}.}
\setcounter{enumi}{10}
\tightlist
\item
  在本小节及下一小节中，我们将考虑一个非空乘积 \(\prod A_{i}\) ，其中 \((A_{i})\) 是E的任意一族（非空）子集。
\end{enumerate}

让 \(J\) 为 \(I\) 的子集. 映射 \((x_{i})_{i\in I}\to(x_{i})_{i\in J}\) （从 \(\prod_{i\in I}A_{i}\) 到 \(\prod_{i\in J}A_{i}\) 上）被称为 \(\prod_{i\in I}A_{i}\) 到 \(\prod_{i\in J}A_{i}\) 上的坐标函数（或投影），并记作 \(pr_{J}\) . 特别地，映射 \((x_{i})_{i\in I}\to x_{x}\) （从 \(\prod_{i\in I}A_{i}\) 到 \(A_{x}\) 上）被称为指标为 \(x\) 的坐标函数（或投影），并记作 \(pr_{x}\) .

如果 \(z\) 是 \(\prod_{i} A_{i}\) 的一个元素，我们有 \(z = (\mathrm{pr}_{i}z)_{i \in I}\) .

设 \(J_1, J_2\) 是构成 I 的一个划分的两个集合。则 \(z \to (\mathrm{pr}_{J_1}z, \mathrm{pr}_{J_2}z)\) 是 \(\prod_{i\in I} A_i\) 到 \(\prod_{i\in J_1} A_i \times \prod_{i\in J_2} A_i\) 上的一个一一映射.

一般来说，如果 \((\mathbf{J}_{\lambda})_{\lambda \in \mathbf{L}}\) 是集合 I 的任意一个划分，映射 \(z \to (\mathrm{pr}_{\mathbf{J}_{\lambda}} z)_{\lambda \in \mathbf{L}}\) 是 \(\prod_{i \in I} A_i\) 到乘积 \(\prod_{\lambda \in L} \left( \prod_{i \in J_{\lambda}} A_i \right)\) 上的一个双射（称为典范的）。这也可以表述为：一族集合的乘积是结合的。

\begin{enumerate}
\def\labelenumi{\arabic{enumi}.}
\setcounter{enumi}{11}
\tightlist
\item
  以下命题推广了§3第3号中的命题； \((\mathbf{X}_t)\) , \((\mathbf{Y}_t)\) 表示E的子集族，使得 \(\mathbf{X}_t \subset \mathbf{A}_t\) 并且 \(\mathbf{Y}_t \subset \mathbf{A}_t\) 对于所有 \(t \in I\) ；Z 表示 的任意子集 \(\prod A_t\) .
\end{enumerate}

\begin{enumerate}
\def\labelenumi{(\alph{enumi})}
\item
  如果 \(\prod_{i} X_{i} \neq \emptyset\) ，关系“ \(\prod_{i} X_{i} \subset \prod_{i} Y_{i}\) “”等价于“对所有 \(i \in I\) , \(X_{i} \subset Y_{i}\) ''.
\item
  我们有 \(\operatorname{pr}_{\chi}^{-1}(\mathbf{X}_{\chi}) = \prod_{i} \mathbf{Y}_{i}\) ，其中 \(\mathbf{Y}_{\chi} = \mathbf{X}_{\chi}\) 并且 \(\mathbf{Y}_{i} = \mathbf{A}_{i}\) 每当 \(i \neq \chi\) 。因此
\end{enumerate}

\[
\prod_ {i \in I} X _ {i} = \bigcap_ {i \in I} \operatorname{pr}_{i}^{-1}(X _ {i}).\tag{47}
\]

\begin{enumerate}
\def\labelenumi{(\alph{enumi})}
\setcounter{enumi}{2}
\tightlist
\item
  如果 \(\prod_{i} X_i \neq \emptyset\) ，那么
\end{enumerate}

\[
\operatorname{pr} _ {x} \left(\prod_ {i} X _ {i}\right) = X _ {x}.\tag{48}
\]

\begin{enumerate}
\def\labelenumi{(\alph{enumi})}
\setcounter{enumi}{3}
\tightlist
\item
  对所有 Z，我们有
\end{enumerate}

\[
\mathbf {Z} \subset \prod_ {\iota} \operatorname{pr} _ {\iota} (\mathbf {Z}).\tag{49}
\]

\begin{enumerate}
\def\labelenumi{(\alph{enumi})}
\setcounter{enumi}{4}
\tightlist
\item
  令 \((\mathbf{J}_1, \mathbf{J}_2)\) 是把 I 划分为两个集合的一个划分，令 \((a_t)_{t \in J_1}\) 是 E 的一族元素，并设 \((\mathbf{X}_t)_{t \in J_2}\) 是 E 的一族子集，使得 \(a_t \in \mathbf{A}_t\) （对所有 \(t \in J_1\) ），以及 \(\mathbf{X}_t \subset \mathbf{A}_t\) （对所有 \(t \in J_2\) ）。则乘积 \(\prod_{i \in I} \mathbf{Y}_t\) ——其中 \(\mathbf{Y}_t = \{a_t\}\) 当 \(t \in J_1\) ，而 \(\mathbf{Y}_t = \mathbf{X}_t\) 当 \(t \in J_2\) ——可以与 \(\prod_{i \in J_2} \mathbf{X}_t\) 通过投影到后一集合上建立一一对应。
\end{enumerate}

\begin{enumerate}
\def\labelenumi{\arabic{enumi}.}
\setcounter{enumi}{12}
\tightlist
\item
  让 \((\mathbf{A}_i)_{i\in I}\) 是某个集合 \(F\) 的一族子集，并且让 \(f\) 是从集合 \(E\) 到乘积 \(\prod_{i\in I} A_i\) 中的一个映射。如果我们令 \(f_i(x) = \text{pr}_i(f(x))\) ，那么 \(f_i\) 是从 \(E\) 到 \(A_i\) 的映射，以及 \(f\) 是映射 \(x \to (f_i(x))\) 。反之，若对每个指标 \(i \in I\) , \(f_i\) 是从 \(E\) 到 \(A_i\) 的映射，那么
\end{enumerate}

\[
x \rightarrow (f _ {i} (x))
\]

是E到 \(\prod_{i\in I}A_i\) 的映射，记为 \((f_i)\) （此处用语略有滥用，因为该记号已表示映射族 \(f_i\) ）。于是我们定义了集合 \(\left(\prod_{i\in I}A_i\right)^E\) 到集合 \(\prod_{i\in I}(A_i^E)\) 的一个双射（称为典范的）.

\begin{enumerate}
\def\labelenumi{\arabic{enumi}.}
\setcounter{enumi}{13}
\item
  设E、F、G为三个集合。对于每个从 \(F \times G\) 到 E 中的映射f，并且对于每个 \(y \in G\) ，令 \(f_{y}\) 表示从F到E的部分映射 \(x \to f(x, y)\) 。则 \(y \to f_{y}\) 是 G 到 \(E^{F}\) 的映射。反之，对于 G 到 \(E^{F}\) 的每个映射 g，存在唯一的从 \(F \times G\) 到 E 的映射 f，使得 \(f_{y} = g(y)\) （对所有 \(y \in G\) ）。因此我们定义集合 \(E^{F \times G}\) 到集合 \((E^{F})^{G}\) 上的一个双射（称为典范双射）.
\item
  让 \((\mathbf{A}_t)_{t \in \mathbf{I}}\) 是集合 \(\mathbf{E}\) 的一族非空子集。对于每一个 \(t \in I\) 让 \(f_t\) 是从 \(\mathbf{A}_t\) 到一个集合 \(\mathbf{F}\) 的映射，使得对于每一对指标 \((t, x)\) , \(f_t\) 和 \(f_x\) 在 \(\mathbf{A}_t \cap \mathbf{A}_x\) 上一致。那么，若
\end{enumerate}

\[
\mathrm{A} = \bigcup_ {i \in \mathrm{I}} \mathrm{A} _ {i},
\]

存在唯一的映射 \(f\) （从 \(\mathbf{A}\) 到 \(\mathbf{F}\) ）使得 \(f\) 到每个 \(\mathbf{A}_i\) 上的限制等于 \(f_i\) 。特别地，如果 \(\mathbf{A}_i \cap \mathbf{A}_x = \emptyset\) 每当 \(i \neq x\) ，我们看到集合 \(\mathbf{F}^{\mathbf{A}}\) 并且 \(\prod_{i \in I} \mathbf{F}^{\mathbf{A}_i}\) 之间存在一一对应（称为典范对应）。

